\chapter{Testes e Avaliação}
\label{cap:test}

Este capítulo apresenta os testes realizados para verificar que a plataforma implementada cumpre os objetivos definidos no Capítulo \ref{chap:necessidades}. Para cada serviço, descreve-se o método de teste, o resultado esperado e o resultado obtido. Complementarmente, discutem-se os resultados de forma crítica, identificando o que funcionou conforme o esperado, as decisões técnicas adotadas e as reflexões sobre o trabalho desenvolvido.

\section{Metodologia de Testes}

Os testes foram estruturados em dois níveis complementares:

\begin{itemize}
    \item \textbf{Testes de funcionamento individual (\textit{smoke tests}):} verificação de que cada serviço arranca corretamente, responde nas portas esperadas e não emite falhas graves (\textit{kernel panics} ou \textit{exits}) nos logs de inicialização.
    \item \textbf{Testes de integração:} verificação de que os serviços interagem corretamente entre si (por exemplo, que a VPN utiliza o Pi-hole como DNS, que o Promtail captura os eventos do Nginx e que a pipeline injeta as credenciais do Vault no container correto).
\end{itemize}

\section{Teste 1: Resolução DNS e Filtragem Ativa (Pi-hole)}

\textbf{Descrição:} Verificar se o Pi-hole responde a pedidos DNS na porta 53 e resolve corretamente nomes de domínio externos e internos.

\textbf{Resultado esperado:} Resolução bem-sucedida de domínios públicos (e.g., \texttt{google.com}), bloqueio de domínios em \textit{blacklists} e latência de resposta otimizada devido à \textit{cache} local.

\textbf{Procedimento:}
\begin{lstlisting}[language=bash, frame=single, breaklines=true]
# Teste de resolucao DNS a partir de um dispositivo externo na LAN
nslookup google.com 192.168.1.252
dig @192.168.1.252 google.com

# Inspecao do comportamento do container
docker logs dns_pihole --tail 20
\end{lstlisting}

\begin{figure}[H]
    \centering
    \includegraphics[width=0.9\textwidth]{pihole_dns_queries.png}
    \caption{Monitorização de consultas DNS bloqueadas e processadas no painel Pi-hole.}
    \label{fig:pihole_dns_queries}
\end{figure}

\textbf{Resultado obtido:} O Pi-hole processou os pedidos corretamente, devolvendo os endereços IP correspondentes sem constrangimentos de \textit{timeout}. A interface web (\texttt{http://192.168.1.252:8080/admin}) validou graficamente as consultas processadas. 

\textbf{Discussão:} O único obstáculo durante a parametrização foi o bloqueio da variável da \textit{password} pela persistência dos volumes (analisado na Secção \ref{sec:pihole}). Ultrapassado este entrave, o serviço manteve-se altamente responsivo. Como trabalho futuro, propõe-se a injeção de registos \textit{DNS A} internos (ex: \texttt{nextcloud.local}) para navegação corporativa simplificada.

\section{Teste 2: Reverse Proxy e Isolamento Read-Only (Nginx)}

\textbf{Descrição:} Verificar se a página HTML institucional é servida corretamente pelo Nginx e testar a inviolabilidade do sistema de ficheiros imposta pela \textit{flag} \texttt{:ro}.

\textbf{Resultado esperado:} A página de estado do *data center* deve carregar normalmente, mas qualquer tentativa interna de modificar o código HTML a partir do processo Nginx deve ser bloqueada ao nível do \textit{kernel}.

\textbf{Procedimento (Verificação HTTP e Read-Only):}
\begin{lstlisting}[language=bash, frame=single, breaklines=true]
# Validar resposta HTTP 200 OK
curl -I http://192.168.1.252

# Tentar violar a seguranca escrevendo no volume
docker exec nginx_proxy touch /usr/share/nginx/html/exploit.txt
\end{lstlisting}

\textbf{Resultado obtido:} O Nginx serviu o site estático corretamente. Ao executar o comando de escrita forçada, o container devolveu o erro \texttt{touch: cannot touch '...': Read-only file system}.

\textbf{Discussão:} O bloqueio de escrita corrobora que o princípio do mínimo privilégio foi aplicado com sucesso. O Nginx atua como um verdadeiro processo descartável, não comprometendo os ficheiros do *host* hospedeiro sob nenhuma circunstância.

\section{Teste 3: Acesso Remoto VPN e Dynamic DNS (WireGuard)}

\textbf{Descrição:} Verificar se um cliente remoto numa ligação de dados móveis (4G) consegue autenticar-se na infraestrutura utilizando o DNS Dinâmico.

\textbf{Resultado esperado:} O túnel estabelece ligação cifrada, resolvendo o domínio do DDNS independentemente do IP do router, garantindo tráfego encapsulado.

\textbf{Procedimento e Resultado:} Após a correção cirúrgica da DMZ do router da operadora (conforme detalhado na Secção \ref{sec:wireguard_debug}), a VPN iniciou o *handshake* de forma instantânea. 
Com o telemóvel desligado do Wi-Fi e conectado na VPN, observou-se:
\begin{itemize}
    \item A pesquisa por "meu IP" na Internet retornou o endereço IP público do servidor residencial, confirmando que o \textit{tunneling} estava ativo de ponta a ponta (\textit{full tunnel}).
    \item Os painéis internos (\texttt{192.168.1.252:3000}) abriram instantaneamente, demonstrando latência negligenciável e acesso seguro às \textit{dashboards} do Grafana fora da LAN.
\end{itemize}

\section{Teste 4: Orquestração CI/CD e Proteção Vault}

\textbf{Descrição:} Verificar se o \textit{GitLab Runner} consegue interceptar novos \textit{commits}, descriptografar os segredos do Ansible Vault e executar o deploy sem interação humana.

\textbf{Resultado esperado:} O estado da pipeline deve reportar \textit{Passed/Succeeded} no repositório web após um \texttt{git push}. O Ansible não deve falhar na desencriptação.

\textbf{Procedimento:} Foram submetidas atualizações à página institucional Web. Verificou-se que a pipeline utilizou a variável \texttt{\$ANSIBLE\_VAULT\_PASSWORD} injetando-a num ficheiro \textit{dummy} transitório, conseguindo processar as passwords das bases de dados, reconfigurar os serviços e apagar o rasto criptográfico, terminando o ciclo com \texttt{failed=0}.

\section{Teste 5: Nuvem Corporativa e E-mail (Nextcloud e Mailserver)}

\textbf{Descrição:} Validar a implementação dos serviços de armazenamento isolado de dados e das comunicações por correio eletrónico.

\textbf{Procedimento e Resultado:} 
O painel do Nextcloud abriu sem falhas críticas, montando os volumes do lado do host de forma permanente, comprovando a eficácia da solução como repositório documental seguro. Relativamente ao e-mail, utilizou-se o comando \texttt{docker exec mailserver setup email add daniel@local.lan}. O processo executou com sucesso e ativou o controlo de quotas Dovecot, validando o aprovisionamento da arquitetura de comunicação corporativa estipulada nos requisitos.

\section{Teste 6: Centralização de Logs e Métricas (Loki, Prometheus e Grafana)}

\textbf{Descrição:} Avaliar a integração da telemetria de hardware e a captura centralizada de eventos estruturados no Grafana.

\textbf{Resultado esperado:} O Loki deve extrair e aplicar *labels* aos logs do Docker, e o Node Exporter deve enviar dados de CPU/RAM para o Prometheus.

\textbf{Procedimento e Resultado:} Interrogou-se a API do Prometheus (\texttt{curl http://192.168.1.252:9090/api/v1/targets}), confirmando que todos os *endpoints* acusavam o estado \texttt{"health":"up"}. Na interface Explore do Grafana, executou-se a *query* LogQL \texttt{\{container="nginx\_proxy"\}}, onde se constataram milhares de linhas de logs perfeitamente indexadas pelo Promtail através da partilha de \textit{socket} do Docker, cumprindo a validação exigida para a observabilidade centralizada.

\section{Análise Crítica dos Resultados}

\subsection{Balanço de Objetivos Cumpridos}

O desenvolvimento demonstrou um saldo muito positivo. Numa análise global, \textbf{100\% dos requisitos delineados no guião foram implementados, testados e validados}:
\begin{itemize}
    \item A transição da infraestrutura física para \textit{containers} Docker Compose resultou numa rede reprodutível em minutos.
    \item O Ansible e o CI/CD (GitLab) eliminaram os pontos cegos de configuração manual, consolidando uma cultura rigorosa de \textit{Infrastructure as Code} (IaC).
    \item As vulnerabilidades de exposição foram resolvidas pela criptografia \textbf{Ansible Vault}, segmentação da rede interna do Docker, acesso VPN através de \textit{Dynamic DNS} e blindagem de volumes.
    \item O ecossistema de serviços de e-mail (Mailserver), DNS (Pi-hole), Reverse Proxy (Nginx) e cloud privada (Nextcloud) funciona de forma fluida.
\end{itemize}

\subsection{O que Poderia ter sido Feito de Forma Diferente (Reflexão)}

Embora o ecossistema atual cubra exaustivamente as necessidades de um laboratório ou PME, o exercício permitiu detetar oportunidades arquiteturais que poderiam ter sido otimizadas noutro contexto:
\begin{itemize}
    \item \textbf{Cifragem HTTPS Real:} Atualmente o Nginx atua sob a porta HTTP 80 internamente. A substituição pelo contentor **Traefik**, ou a automatização da emissão de certificados \textit{Let's Encrypt} com a ferramenta **Certbot**, seria um avanço orgânico essencial se a infraestrutura precisasse de expor o \textit{Nextcloud} ou serviços web ao tráfego público de forma direta na porta HTTPS 443 \cite{rfc8446}.
    \item \textbf{Gestão Ativa de Alertas:} O Prometheus e o Grafana oferecem painéis analíticos visualmente sofisticados para \textit{troubleshooting} reativo, contudo, a plataforma carece de \textit{trigger rules}. A configuração do \textit{Alertmanager} permitiria o envio proativo de mensagens (ex: Telegram ou E-mail corporativo) sempre que a VM atingisse picos críticos de consumo de CPU/RAM.
\end{itemize}

Conclui-se que o rigor na documentação de incidentes (ex: o problema de volume no Pi-hole, o socket do Promtail e o troubleshooting da DMZ da MEO) moldou o aspeto central deste trabalho prático, traduzindo as falhas momentâneas em valiosas referências pedagógicas no domínio da Engenharia de Redes.